% Created 2026-05-22 Fri 11:10
% Intended LaTeX compiler: xelatex
\documentclass[11pt]{article}
\usepackage{graphicx}
\usepackage{longtable}
\usepackage{wrapfig}
\usepackage{rotating}
\usepackage[normalem]{ulem}
\usepackage{capt-of}
\usepackage{hyperref}
\usepackage{amsmath}
\usepackage{amssymb}
\usepackage{graphicx}
\usepackage{geometry}
\usepackage[a4paper,margin=0.75in,top=0.85in]{geometry}
\usepackage{microtype}
\usepackage{parskip}
\usepackage{titlesec}
\titleformat{\section}{\large\bfseries}{}{0em}{}[\titlerule]
\titleformat{\subsection}{\normalize\bfseries}{}{0em}{}
\usepackage{enumitem}
\setlist{noitemsep, topset=4pt}
\usepackage{titling}
\pretitle{\begin{center}\large\bfseries}
\posttitle{\end{center}}
\author{Tyrone}
\date{\today}
\title{ENGDATA Notes}
\hypersetup{
 pdfauthor={Tyrone},
 pdftitle={ENGDATA Notes},
 pdfkeywords={},
 pdfsubject={},
 pdfcreator={Emacs 30.2 (Org mode 9.7.11)}, 
 pdflang={English}}
\begin{document}

\maketitle
\tableofcontents


\noindent\rule{\textwidth}{0.5pt}
\section{Type of Data}
\label{sec:org1e67d3c}

Study of the population is a \textbf{census}. Study of a sample of a population is a \textbf{survey}.

Data can either be numerical or categorical. Just because the data points are numbers, it doesn't mean they are numerical (e.g. student ID numbers, years).

The \textbf{Delphi technique} is when you send rounds of questionnaires to a panel of experts to reach a consensus.
\section{Types of Probability Sampling}
\label{sec:orga6cff89}

\begin{enumerate}
\item \uline{Simple Random Sampling}: Every member has an equal chance of being selected (e.g. lottery, random number generator)

\item \uline{Cluster Sampling}: The population is divided by clusters, which is often geographical. And entire clusters are randomly selected.

\item \uline{Stratified Sampling}: The population is divided into subgroups (strata) based on shared characteristics and random samples are taken from each strata.

\item \uline{Systematic Sampling}: Every n-th member of a group is selected.
\end{enumerate}


The point of probability sampling is that each individual has a known, non-zero chance of being selected; minimizing bias and allowing for generalization.
\section{Types of Non-Probability Sampling}
\label{sec:org915d821}

\begin{enumerate}
\item \uline{Convenience Sampling}: Individuals are selected for their availability and accessibility.

\item \uline{Voluntary Respone Sampling}: Individuals select themselves.

\item \uline{Purposive/Judgmental Sampling}: The researcher uses their expertise to select a sample that is most useful to the purpose of the research.

\item \uline{Snowball Sampling}: Participants refer other participants.

\item \uline{Quota Sampling}: Selecting a specific number of individuals from different subgroups.
\end{enumerate}

These methods are faster and easier but risks higher bias. Most often used in exploratory research.
\section{Types of Comparisons}
\label{sec:org8248ffb}

\begin{enumerate}
\item \uline{Component comparison}: "How does a part relate to the whole?" (pie chart)

\item \uline{Correlation comparison}: "Do two variables move together?" (scatter plot, bubble chart)

\item \uline{Frequency comparison}: "How many items fall into each range?" (histogram, bell curve)

\item \uline{Item comparison}: "How do things rank against each other?" (bar chart, column chart)

\item \uline{Time series}: "How does something change over time?" (line chart, area chart)

\noindent\rule{\textwidth}{0.5pt}
\end{enumerate}
\section{Probability}
\label{sec:org65d74b6}

\subsection{\textbf{Defining Sample Space}}
\label{sec:org3da2178}

\uline{Listing/Roster method}: enumerating all elements of the sample space.

\uline{Property/Set-Builder method}: Choosing a property or characteristic common to all sample points and then using this common characteristic to define the sample space.

\uline{Tree Diagram}: Uses a tree data structure to express sequential order of events of an experiment in chronological order.

An \textbf{event} is a subset of a sample space.

A \textbf{sample point} is a specific outcome from the \textbf{sample space}.

A \textbf{simple event} is an event which only has one outcome.

A \textbf{compound event} is an event which consists of more than one outcome and may be decomposed into simple events.
\subsection{\textbf{Operations on Events}}
\label{sec:org3e50bd5}

\uline{Union}: A U B, set of all events in A or B.

\uline{Intersection}: A \(\cap\) B. Set of all events in A and B

\uline{Complement}: A' represents the set of all events not found in A.

\textbf{S} represents the sample space. So S' represents all events not in the sample space.

A \& B = NULL is \textbf{mutually exclusive} or \textbf{disjointed}.

\noindent\rule{\textwidth}{0.5pt}
\section{Counting Principles and Techniques}
\label{sec:org8e387f3}

\subsection{Counting Theory}
\label{sec:org419d6b0}
A set of techniques used for identifying the total number of possibilities of a statistical experiment without enumerating the elements of the sample space.

\uline{Multiplication Rule}:

n\textsubscript{T} = n\textsubscript{1} \texttimes{} n\textsubscript{2} \texttimes{} \ldots{} \texttimes{} n\textsubscript{k}

Where \emph{n} is an outcome.

\uline{Addition Rule}:
If \emph{k} outcomes are mutually exclusive, and outcome 1 can be performed in m\textsubscript{1} ways, outcome 2 can be performed in m\textsubscript{2} ways, and so on for \emph{k} outcomes, then the total number of outcomes is given by:

m\textsubscript{T} = m\textsubscript{1} + m\textsubscript{2} + \ldots{} + m\textsubscript{k}
\subsection{Simple Permutation}
\label{sec:orga6a37f7}

There is an \uline{ordered} arrangement of \textbf{distinct objects}.

The formula:

\[
nPr = \frac{n!}{(n-r)!}
\]

Where \emph{n} is the total number of objects and \emph{r} is how many are being selected from \emph{n}.
\subsection{Simple Combination}
\label{sec:org319097d}

An arrangement of distinct objects \uline{without regard to order}.

\[
nCr = \frac{n!}{(n-r)! * r!}
\]
\subsection{Special types of permutation}
\label{sec:org1f6e5a7}

\subsubsection{Circular Permutation}
\label{sec:orga20f324}

Arrangement of objects in a circle.

The formula is:

\[
Pnc = (n-1)!
\]
\subsubsection{Arrangement of Similar Objects}
\label{sec:orge8ccc9a}

To find the number of permutations of \emph{n} "objects" \uline{all taken at the same time}.

\begin{itemize}
\item n\textsubscript{1} are objects of type 1.
\item n\textsubscript{2} are objects of type 2.
\item n\textsubscript{k} are objects of type \emph{k}.
\end{itemize}

\emph{- n} represents the total objects.


\[
\sum\limits_{i=1}^{k}^{} n_{i}_{} = n
\]

This translates to:

\[
\frac{n!}{n_{1}! n_{2}! n_{3}! ... n_{k}!}
\]
\subsubsection{Partitioning of n objects into cells}
\label{sec:org0e77a10}

To find the number of permutations of \emph{n} "objects" \uline{all taken at the same time}.

\begin{itemize}
\item n\textsubscript{1} are the objects placed in cell 1.
\item n\textsubscript{2} are the objects placed in cell 2.
\item n\textsubscript{k} are the objects placed in cell \emph{k}.
\end{itemize}

Very similar to the previous type of permutation.


\[
\sum\limits_{i=1}^{k}^{} n_{i}_{} = n
\]

This translates to:

\[
\frac{n!}{n_{1}! n_{2}! n_{3}! ... n_{k}!}
\]
\section{Axioms of Probability}
\label{sec:org6b2f50a}

\uline{Axiom 1}: 0 <= P(A) <= 1.

\uline{Axiom 2}: P(A) = 1 if A is a certain or sure event.

\uline{Axiom 3}: P(A) = 0 if A is an impossible event.

\uline{Axiom 4}: If the probability of events in the sample space is added, the sum is always 1.

\[
\sum\limits_{i=1}^{n} P(E) = 1
\] 
\section{Ways of measuring probability}
\label{sec:org407d619}

\subsection{Classical Concept}
\label{sec:orgd1afeb8}

The probabilities of events can be determined without gathering data.

\[
P(S) = \frac{n(A)}{N}
\]

Where \emph{n(A)} is the subset of the sample space and \emph{N} is the total sample space.
\subsection{Relative Frequency}
\label{sec:orgf0e611a}

Uses historical data or survey results to determine probabilities.

\[
P(A) = \frac{f_{A}}{f_{T}}
\]

Where f\textsubscript{A} is the frequency of occurance of event A and f\textsubscript{T} is the total number of observations.
\subsection{Subjective Assessment}
\label{sec:orgcb2d69d}

An expert in the field provides an educated guess of the probability.
\section{Probability Rules}
\label{sec:org8283439}

\subsection{Union (A or B)}
\label{sec:org206559d}

\[
P(A U B) = P(A) + P(B) - P(A \cap B)
\]

If A and B are mutually exclusive: \(P(A U B) = P(A) + P(B)\) 
\subsection{Intersection (A and B)}
\label{sec:orgf551cbe}

\[
P(A \cap B) = P(A) \times P(B | A)
\]

If A and B are independent: \(P(A \cap B) = P(A) \times P(B)\)
\subsection{Complement of A}
\label{sec:org1437a2e}

\[
P(A') = 1 - P(A)
\]
\subsection{Probability of B given A}
\label{sec:org3e58e1d}

\[
P(B|A) = \frac{P(A \cap B)}{P(A)}
\]

\begin{enumerate}
\item Find the probablity of A and B happening.

\item Divide the probability of both events happening by the probability of event A
\end{enumerate}

If A and B are independent:


$$\begin{align}
P(B|A) &= \frac{P(A \cap B)}{P(A)} \tag{1}\\
P(B) &= \frac{P(A \cap B)}{P(A)} \tag{2}\\
P(B) \cdot P(A) &= P(A \cap B) \tag{3}\\
P(A \cap B) &= P(A) \cdot P(B) \tag{4} 
\end{align}$$
\subsection{Venn Diagram}
\label{sec:org38ed9bb}

The entire sample space can be expressed as:

\[
P(Sample Space) = P(A \cap B') + P(A' \cap B) + P(A \cap B)  + P(A' \cap B')
\]
\subsection{Contingency Table}
\label{sec:orga92e4a9}

A tool used to facilitate the calculation of probabilities involving two events.

\begin{table}[h]
\centering
\begin{tabular}{l c c c}
\hline
         & $B$              & $B'$               & Total   \\
\hline
$A$      & $P(A \cap B)$    & $P(A \cap B')$     & $P(A)$  \\
$A'$     & $P(A' \cap B)$   & $P(A' \cap B')$    & $P(A')$ \\
\hline
Total    & $P(B)$           & $P(B')$            & $1$     \\
\hline
\end{tabular}
\end{table}
\subsection{Conditional Probability}
\label{sec:org9ae99d4}

The probability that an event will occur given that some other event has already occured.

\[
P(A|B) = \frac{P(A\capB)}{P(B)}
\]
\subsection{Binomial Probability}
\label{sec:orgd25d6a0}
\section{Bayes Theorem}
\label{sec:orga142557}

Unlike usualky where P(effect|cause) is calculated, Bayesian reasoning tries to find P(cause|effect).

\[
P(A | B) = \frac{P(B | A) \times P(A)}{P(B)}
\]
\end{document}
